ट्यूरिंग यंत्रे त्यांच्या सूचनासंचांनी निर्धारित होतात. हे सूचनासंच पंचकांचे सांत संच असतात (प्रत्येक अवस्था आणि वाचलेले चिन्ह यांसाठी नवी अवस्था, लिहायचे चिन्ह आणि शीर्षाची हालचाल दिलेली असते). पंचकांचे सांत संच गणनीय आहेत, म्हणून ट्यूरिंग यंत्राच्या प्रत्येक सूचनासंचाशी एक संख्या जोडण्याची पद्धत आहे. अशा एका निश्चित पद्धतीने एखाद्या ट्यूरिंग यंत्राच्या सूचनासंचाशी जोडलेली संख्या म्हणजे त्या यंत्राचा \emph{निर्देशांक}. अशा रीतीने ट्यूरिंग यंत्रांविषयी अप्रत्यक्षपणे, म्हणजे त्यांच्या निर्देशांकांविषयी बोलून, ``बोलता येते''.

ट्यूरिंग यंत्रांच्या वर्तनाविषयी एक महत्त्वाची समस्या म्हणजे ती अखेरीस थांबतात की नाही हे ठरवणे. $h(e, n)$ असे फलन असू द्या की निर्देशांक~$e$ असलेले ट्यूरिंग यंत्र आदान~$n$ वर सुरू केले असता थांबते, तर त्याचे मूल्य $= 1$ आहे, आणि अन्यथा $=0$ आहे. त्याला \emph{थांबण्याचे फलन} म्हणतात. थांबण्याचे फलन स्वतः ट्यूरिंग-संगणनक्षम आहे का, या प्रश्नाला \emph{थांबण्याची समस्या} म्हणतात. उत्तर नकारार्थी आहे. थांबण्याची समस्या सोडवता येत नाही. हे कर्णरेषीय युक्तिवादाने सिद्ध केले जाते.

``ट्यूरिंग यंत्रे वापरून $X$ साधता येईल का?'' या स्वरूपाच्या समस्यांच्या एका मोठ्या वर्गाचे थांबण्याची समस्या हे केवळ एक उदाहरण आहे. तर्कशास्त्राची आणखी एक मध्यवर्ती समस्या म्हणजे \emph{प्रथम-क्रम तर्कशास्त्राची निर्णय समस्या}. दिलेले वाक्य वैध आहे की नाही हे ठरवू शकेल असे ट्यूरिंग यंत्र आहे का, असा तो प्रश्न आहे. या प्रसिद्ध समस्येचेही उत्तर नकारार्थी आहे. निर्णय समस्या सोडवता येत नाही. हे न्यूनीकरणाच्या युक्तिवादाने सिद्ध केले जाते. प्रत्येक ट्यूरिंग यंत्र~$M$ आणि आदान~$w$ यांच्याशी असे एक प्रथम-क्रम वाक्य~$T(M, w) \lif E(M, w)$ जोडता येते की $M$ हे आदान~$w$ वर सुरू केले असता थांबते तेव्हा आणि केवळ तेव्हाच ते वाक्य वैध असते. निर्णय समस्या सोडवता आली असती, तर तिचा वापर करून थांबण्याची समस्याही सोडवता आली असती.
